% ============================================================================
% Fichier  : partie5/P05_C14_enquete_integratrice.tex
% Titre    : L'Enquête Intégratrice
% Auteur   : MINKA MI NGUIDJOÏ Thierry Emmanuel
% Version  : 1.0 -- novembre 2026
% Corrigés : corriges/Corr_C14.tex
% ============================================================================
% Chapitre de synthèse : reprend les Chapitres 1 à 13 dans une méthode en dix
% étapes, appliquée à NYANGA (récapitulatif) puis proposée pour SOLITRANS
% (février 2026, exercice intégrateur et lab FAV-LAB-10).
% Sources : ISO/IEC 27035, 27043 ; NIST SP 800-61 r2 et 800-86.
% Faits nouveaux (bible) : coût de l'incident NYANGA 118 M FCFA (47,3 M de
% virements dont 12,6 M recouvrés, enquête 21 M, remplacement et mesures
% 34 M, arrêt des agences 28,3 M) ; durée de la mission SANKOFA 23 jours.
% ============================================================================

\chapter{L'Enquête Intégratrice}
\label{chap:enquete}
\chaptermark{L'enquête intégratrice}

\epigraph{%
    « Une enquête bien conduite se reconnaît à ce qu'un autre\\
    pourrait la refaire à partir du dossier. »%
}{Principe du cours}

\niveauChapitre{Expert}{%
    Maîtrise des Chapitres~\ref{chap:investigation} à~\ref{chap:antiforensique}
}

\begin{boxobjectifs}
\begin{itemize}
    \item Conduire une enquête complète en dix étapes, de l'alerte à la
          clôture.
    \item Articuler réponse à incident, enquête technique et enquête
          judiciaire sur une même affaire.
    \item Tenir le dossier d'enquête de façon qu'il soit reproductible.
    \item Relier l'enquête au programme de gestion des vulnérabilités.
    \item Estimer le coût d'un incident et le présenter à une direction.
\end{itemize}
\end{boxobjectifs}

\bigskip

% ============================================================================
\section*{Avant de commencer : vingt-trois jours}
% ============================================================================

Yaoundé, 6 avril 2025. La mission de SANKOFA EXPERTISE à NYANGA CRÉDIT
s'achève après vingt-trois jours. Mme~MBARGA réunit son équipe pour le
dernier exercice de chaque mission~: relire le dossier comme le relirait un
expert adverse, page par page, en se demandant pour chaque conclusion d'où
elle vient. Elle ne cherche pas les erreurs de l'attaquant, mais les siennes.

\begin{boxalert}
\textbf{Violation typique.} Conduire chaque étape correctement, mais sans
fil~: des acquisitions sans questions, des analyses sans chronologie, une
chronologie sans arbre des causes, un rapport sans recommandations
vérifiables.

\smallskip
\textbf{Règle~:} une enquête suit une méthode écrite, dont chaque étape
produit une pièce du dossier et prépare la suivante. Le dossier, pas la
mémoire des enquêteurs, porte l'enquête.
\end{boxalert}

\separateur

% ============================================================================
\section{Une Méthode en Dix Étapes}
\label{sec:c14-methode}
% ============================================================================

\begin{tableaucours}{cL{3.4cm}L{5.0cm}L{2.6cm}}{%
    \tableauHeader{N°} &
    \tableauHeader{Étape} &
    \tableauHeader{Pièce produite} &
    \tableauHeader{Chapitre}
}
1 & Recevoir l'alerte, isoler & Registre des actions ouvert & \ref{chap:prologue} \\
2 & Cadrer~: mandat, questions, cadre juridique & Lettre de mission,
questions écrites & \ref{chap:investigation}, \ref{chap:normes-droit} \\
3 & Formuler les hypothèses & Liste des hypothèses et pièces réfutantes &
\ref{chap:investigation} \\
4 & Préserver et acquérir & Formulaires de chaîne de possession, scellés &
\ref{chap:preuve}, \ref{chap:acquisition} \\
5 & Examiner disque, mémoire, réseau & Notes d'examen, sorties d'outils &
\ref{chap:disque} à \ref{chap:reseau} \\
6 & Construire la chronologie & Tableau des horloges, chronologie cotée &
\ref{chap:chronologie} \\
7 & Remonter à la vulnérabilité & Matrice des explications, arbre des
causes & \ref{chap:artefact-vuln} \\
8 & Chercher les variantes & Liste des autres instances & \ref{chap:artefact-vuln} \\
9 & Rendre compte & Rapport, annexes, recommandations &
\ref{chap:rapport} \\
10 & Clore et transmettre & Restitution ou transfert des scellés, règles
nouvelles au programme & \ref{chap:gerer} \\
\end{tableaucours}
\vspace{-0.8em}
\begin{center}{\small\itshape Tableau~14.1. La méthode du cours en dix étapes}\end{center}

Les étapes se chevauchent~: on acquiert encore pendant qu'on examine, et
chaque découverte peut rouvrir une hypothèse. La méthode fixe l'ordre des
dépendances, pas un calendrier. Elle respecte les classes de processus de
l'ISO/IEC~27043~\cite{iso27043_2015} et les phases du NIST
SP~800-86~\cite{nist80086}, en y ajoutant explicitement les étapes~7 et~8,
qui relient l'enquête à la gestion des vulnérabilités.

% ============================================================================
\section{Trois Enquêtes, un Dossier}
\label{sec:c14-trois}
% ============================================================================

À NYANGA CRÉDIT, la réponse à incident, l'enquête du cabinet et l'enquête
judiciaire ont avancé ensemble. Leur coordination tient à quelques règles.

\begin{tableaucours}{L{4.4cm}L{7.8cm}}{%
    \tableauHeader{Règle} &
    \tableauHeader{Application à NYANGA}
}
Un seul registre des actions & Ouvert par la RSSI le 14~mars, repris par le
cabinet, versé à la procédure \\
Les originaux vont à la justice & Scellés NC-01 à NC-07 remis à l'OPJ le
16 et le 22~mars~; le cabinet travaille sur ses copies \\
Chacun garde son rapport & Le cabinet rend compte à l'établissement,
l'expert au procureur \\
La reprise d'activité attend la préservation & Contrôleur de secours en
service après la capture de la mémoire et l'image du disque \\
Les communications passent par la cellule & Aucune déclaration publique
avant validation par la juriste \\
\end{tableaucours}
\vspace{-0.8em}
\begin{center}{\small\itshape Tableau~14.2. Coordonner les trois enquêtes}\end{center}

\begin{boxinfo}
\textbf{Les obligations de notification.} Un incident qui expose des données
de clients peut déclencher des obligations de notification envers une
autorité de protection des données et envers les personnes concernées, au
titre de la loi n°~2024/017~\cite{loi2024017}. Un établissement de
microfinance peut aussi devoir informer son autorité de supervision [nature
et délai à vérifier dans la réglementation applicable]. Le détail relève de
l'UE MSN-2S3-JUR01~; l'enquêteur fournit les faits qui permettent de
remplir ces obligations~: quelles données, combien de personnes, depuis
quand.
\end{boxinfo}

% ============================================================================
\section{Le Dossier NYANGA}
\label{sec:c14-dossier}
% ============================================================================

\begin{boxscenario}{Opération NYANGA. Le dossier au 6 avril 2025}
\begin{tableaucours}{L{3.8cm}L{8.4cm}}{%
    \tableauHeader{Rubrique} &
    \tableauHeader{Contenu}
}
Mandat et questions & Lettre de mission du 14~mars~; cinq questions (point
d'entrée, séquence, données sorties, auteurs, mesures) \\
Registre des actions & 214 entrées, du 14~mars 8~h~15 au 6~avril \\
Scellés & NC-01 à NC-07, formulaires, procès-verbaux de remise \\
Notes d'examen & Disque (SRV-CPT01, SRV-AD01, SRV-FIC01), mémoire
(SRV-AD01, SRV-FIC01), réseau (NC-04), sorties brutes des outils en
annexe \\
Chronologie & 14 lignes cotées, tableau des horloges, temps de présence \\
Causes & Matrice E1-E2-E3, arbre des sept causes \\
Variantes & Recherche lancée, programme de septembre \\
Rapport & Version du 6~avril, relue contradictoirement \\
\end{tableaucours}

\medskip
\textbf{Réponses aux cinq questions.} Point d'entrée~: \OUI{}
(VULN-2024-4471). Séquence~: \OUI{}, sauf du 3 au 8~mars (\NEK{}). Données
sorties~: \OUI{} pour le dossier \texttt{Clients} par le siège, \NEK{} pour
d'autres chemins. Auteurs~: \NEK{}. Mesures~: sept recommandations.
\end{boxscenario}

\begin{boxlecture}{Les cinq questions de la mission}
\textbf{Décompte.} Trois questions établies (point d'entrée, séquence
principale, données sorties par le siège), une non tranchée (auteurs), et
deux réserves partielles (trou du 3 au 8~mars, autres chemins de sortie).

\medskip
\textbf{Lecture.} La mission répond à ce que l'établissement doit savoir
pour décider~: fermer la porte, mesurer la fuite, prévenir les clients. Elle
ne répond pas à la question pénale de l'identité, qui dépend de pièces
détenues par des tiers à l'étranger.
\end{boxlecture}

% ============================================================================
\section{Le Coût d'un Incident}
\label{sec:c14-cout}
% ============================================================================

La direction demande toujours combien l'incident a coûté. La réponse se
construit poste par poste, en distinguant ce qui est constaté (factures,
pertes comptabilisées) de ce qui est estimé.

\begin{boxscenario}{Opération NYANGA. Le coût au 30 juin 2025}
\begin{tableaucours}{L{6.6cm}rL{2.6cm}}{%
    \tableauHeader{Poste} &
    \tableauHeader{M FCFA} &
    \tableauHeader{Nature}
}
Virements frauduleux & 47,3 & Constaté \\
Montants recouvrés par les opérateurs & $-12{,}6$ & Constaté \\
Mission d'enquête et d'expertise privée & 21,0 & Constaté (factures) \\
Remplacement de la passerelle et mesures immédiates & 34,0 & Constaté \\
Arrêt partiel des agences (estimation des pertes d'activité) & 28,3 &
Estimé \\
\textbf{Total} & \textbf{118,0} & \\
Notification et accompagnement des clients exposés & Non chiffré & \NEK{} \\
Contentieux éventuels & Non chiffré & \NEK{} \\
\end{tableaucours}

\medskip
\textbf{Analyse} (confiance élevée sur les postes constatés, moyenne sur
l'arrêt des agences). Le coût connu, 118~M~FCFA, est une borne basse. Les
deux postes \NEK{} peuvent le dépasser. À titre de comparaison, la mise à
jour de la passerelle en décembre 2024 aurait demandé une intervention de
deux heures du prestataire.
\end{boxscenario}

% ============================================================================
\section{L'Exercice Intégrateur : SOLITRANS, Février 2026}
\label{sec:c14-solitrans}
% ============================================================================

Le portail des transitaires de SOLITRANS a été compromis en février 2026
(Chapitre~\ref{chap:antiforensique}). Le laboratoire FAV-LAB-10 de NDAP
SUKULU en fournit les pièces~: journaux d'audit de la console du
fournisseur, journaux de l'orchestrateur, journaux de l'application envoyés
au collecteur, instantané des disques, historique des balayages du portail,
tickets de correction. Le lecteur y applique les dix étapes. Trois
questions guident l'exercice~:
\begin{enumerate}
    \item Par quelle faille l'attaquant est-il entré, et cette faille
          était-elle connue du programme de gestion des vulnérabilités~?
    \item Qu'a-t-il modifié, et qu'a-t-il pu lire~?
    \item Le programme de SOLITRANS a-t-il échoué, ou l'attaque est-elle
          passée par une faille hors de sa portée~?
\end{enumerate}

\begin{boxinfo}
\textbf{Une question piège.} La troisième question n'a pas de réponse
évidente. Une faille connue, classée, dont l'exception E-03 était en cours,
n'est pas un échec du programme au même titre qu'une faille ignorée~; mais
une mesure compensatoire défaillante en est un. L'exercice demande de
coter, pas de défendre.
\end{boxinfo}

% ============================================================================
\section{Ce que l'Enquête demande à chaque Métier}
\label{sec:c14-metiers}
% ============================================================================

\begin{boxmetier}{Analyste en réponse à incident}
\textbf{Relire son dossier avant de le remettre}
\begin{enumerate}
    \item Chaque conclusion renvoie-t-elle à des pièces cotées~?
    \item Chaque pièce a-t-elle sa chaîne de possession~?
    \item Chaque outil a-t-il sa version et sa sortie brute en annexe~?
    \item Chaque heure a-t-elle son fuseau et sa dérive~?
    \item Chaque \NEK{} est-il dit, avec ce qui permettrait de le lever~?
\end{enumerate}
\end{boxmetier}

\begin{boxmetier}{RSSI}
Présenter à la direction le coût constaté, le coût estimé et ce qui reste
inconnu, puis le coût de la mesure qui aurait évité l'incident. La
comparaison fait plus pour le budget de la sécurité que toute projection.
\end{boxmetier}

\begin{boxmetier}{Analyste en gestion des vulnérabilités}
Recevoir le dossier d'enquête et répondre par écrit à une question~: le
programme aurait-il dû voir cette faille, et pourquoi ne l'a-t-il pas vue~?
\end{boxmetier}

\begin{boxmetier}{Application OPJ/Expert Judiciaire}
L'expert relit son rapport en se posant la question de l'expert adverse~:
quelle conclusion pourrais-je attaquer, et sur quelle pièce~?
\end{boxmetier}

\separateur

\begin{boxresume}
\textbf{Ce qu'il faut retenir de ce chapitre~:}
\begin{itemize}
    \item Une enquête suit dix étapes, chacune produisant une pièce du
          dossier~; les étapes~7 et~8 la relient à la gestion des
          vulnérabilités.
    \item Réponse à incident, enquête privée et enquête judiciaire
          partagent un registre et des scellés, pas un rapport.
    \item Le coût d'un incident se construit poste par poste, constaté ou
          estimé~; ce qui reste \NEK{} se dit.
    \item À NYANGA, trois des cinq questions sont tranchées, l'identité des
          auteurs reste \NEK{}, et le coût connu atteint 118~M~FCFA.
\end{itemize}
\end{boxresume}

\begin{boxinfo}
\textbf{Points de contrôle.} Vous avez compris ce chapitre si vous savez~:
\begin{itemize}
    \item citer les dix étapes et la pièce que chacune produit~;
    \item organiser la coordination de trois enquêtes sur une affaire~;
    \item présenter un coût d'incident en distinguant constaté, estimé et
          inconnu.
\end{itemize}
\end{boxinfo}

% ============================================================================
\section*{Exercices}
% ============================================================================

\begin{boxexo}[title={\ding{45}\quad Exercice~14.1. Ordonner \niveaubase}]
Remettez dans l'ordre de la méthode~: (a) construire la chronologie~;
(b) formuler les hypothèses~; (c) isoler~; (d) chercher les variantes~;
(e) acquérir~; (f) remonter à la vulnérabilité~; (g) cadrer le mandat~;
(h) examiner~; (i) rendre compte~; (j) clore.
\end{boxexo}

\begin{boxexo}[title={\ding{45}\quad Exercice~14.2. Relecture contradictoire \niveaumoyen}]
Jouez l'expert adverse. Choisissez trois conclusions du dossier NYANGA et,
pour chacune, formulez l'objection la plus forte possible, puis la réponse
que le dossier permet d'y apporter. Indiquez si l'une d'elles devrait
changer de cote.
\end{boxexo}

\begin{boxexo}[title={\ding{45}\quad Exercice~14.3. Le coût \niveaumoyen}]
\begin{enumerate}[(a)]
    \item Vérifiez le total de l'encadré « Le coût au 30 juin 2025 ».
    \item Quelle part du coût connu est estimée plutôt que constatée~?
    \item Rédigez trois phrases pour présenter ce coût au conseil
          d'administration.
\end{enumerate}
\end{boxexo}

\begin{boxexo}[title={\ding{45}\quad Exercice~14.4. SOLITRANS, février 2026 \niveauexpert}]
À partir des pièces du laboratoire FAV-LAB-10, conduisez les dix étapes et
répondez aux trois questions de la section~14.5. Pour chaque réponse,
donnez une cote et les pièces qui la fondent. Rédigez enfin la note d'une
page destinée à M.~NJIKE.
\end{boxexo}

\bigskip

\begin{boxinfo}
\textbf{Pour aller plus loin.}
\begin{itemize}
    \item L'ISO/IEC~27043~\cite{iso27043_2015} et l'ISO/IEC~27035~\cite{iso27035}.
    \item Le NIST SP~800-61~\cite{nist80061r2}.
\end{itemize}
\end{boxinfo}

\begin{boxinfo}
\textbf{Transition.} Le dossier est complet. Il reste à l'écrire pour ceux
qui ne l'ont pas vu se construire, et à le défendre devant ceux qui le
contestent~: Chapitre~\ref{chap:rapport}.
\end{boxinfo}
